\section{Related Work}
\label{sec:related}

The comparison separates condition discovery, safe executable checking,
conditional candidate proofs, and complete-program installation. Evidence for
one link does not imply all four, and names such as ``optimistic'' or ``peephole''
do not determine an interface's expressive power.

\paragraph{Optimistic loops and candidate conditioning.}
Doerfert et al.'s OLO~\cite{olo2017} organizes assumptions from polyhedral
modeling, derives and simplifies runtime conditions, and handles conditional
preloads and check arithmetic. These are direct algorithmic precedents for the
case study. Our comparison concerns a mechanical proof chain for a declared
Clight source class, not the invention of runtime loop versioning. OLO's
production compiler and performance results also set a practical standard that
our current proof endpoints alone do not meet.

COVE/cSTOKE~\cite{cove2015} accepts independently produced program candidates,
infers sufficient conditions using tests and analysis, and verifies conditional
equivalence. Its dynamic-check experiments use C instrumentation compiled and
joined with the optimized binary. This makes candidate conditioning a broader
motivation than loop optimization. GuardCert's current template-based adapters
do not provide COVE's general condition-inference capability. Our relevant
comparison is the proved connection between the actual generated check, its
safe evaluation domain, the candidate, and installation.

\paragraph{Verified speculation and recovery.}
CoreJIT~\cite{corejit2021} mechanically verifies speculative optimization and
deoptimization in an IR with explicit assumptions and synchronization points.
It separates optimization policy from verified execution mechanisms. This is a
precedent for verified dynamic assumptions, not merely an equivalence checker.
Our present entry-versioning contract retains an original fragment as fallback;
it does not reconstruct an active frame at an internal deoptimization point.
CoreJIT's 2021 native experiment lies outside that paper's proof boundary.
A later effectful-JIT development~\cite{fmjit2023} connects native generation to
the CompCert backend under its own scope. These two results should not be
combined into an unstated theorem about all speculative passes.

\paragraph{Local validation and whole-program lifting.}
Peek~\cite{peek2016} provides verified assembly peephole rules and a
liveness-based connection from local simulation to complete programs. Its
published framework imposes replacement/control restrictions and source
normalization obligations. These establish a substantial local-to-global
precedent. A comparison must identify the extra check-safety, private-state,
and source-prefix obligations of our loop example; moving the example to
Clight or adding a conditional is insufficient.

Chamois~\cite{chamois2023} validates oracle-proposed control-flow graphs using
block simulations and relational invariants, with integration into CompCert.
Its documented memory-copy expansion proves concrete CFG and private-resource
simulation~\cite{chamoisapi}. We have inspected those interfaces, but have not
instantiated its artifact with our loaded-bound example. The absence of such a
comparison is not evidence that its framework cannot express a guarded rewrite.

\paragraph{Comparison still required.}
\pending{For one fixed alias-sensitive loaded-bound source, record which
obligations and proof services OLO, CoreJIT, Chamois, Peek, and GuardCert supply.
Where artifacts are instantiated, report the additional source/model,
condition-safety, frame, and placement work. Distinguish unknown capabilities
from demonstrated restrictions before making a novelty or author-effort claim.}
